% GENERATED FILE. Edit canonical lessons, then run npm run generate:books.
% canonical-source-hash: fnv1a64:8a3ed6e4049cf951
% canonical-lessons: AR-C526-mikaniki, AR-C526-muqawil, AR-C526-dahhan, AR-C526-banna, AR-C526-iskafi

\chapter{Trades and crafts}
\label{ch:ar-trades-and-crafts}
\hlchaptermodality{526}

\begin{chapteropening}
\textbf{By the end of this chapter:} \emph{I can name a mechanic, a contractor, a house painter, a builder and a cobbler.}

Complete the last lesson of chapter 526: I can name a mechanic, a contractor, a house painter, a builder and a cobbler.
\end{chapteropening}

\section[\texorpdfstring{\ar{ميكانيكي} (mīkānīkī)}{ميكانيكي (mīkānīkī)}]{\ar{ميكانيكي} (mīkānīkī) — a mechanic}

\label{lesson:AR-C526-mikaniki}



\begin{quote}
Before the new one: say the Arabic for to program, then the Arabic for a programmer.
\end{quote}

\subsection*{You'll want to know: \ar{ميكانيكي}}
\textbf{\ar{ميكانيكي}} — \emph{mīkānīkī} — \textquotedblleft{}a mechanic\textquotedblright{}.

\begin{quote}
\ar{أصلح} \ar{الميكانيكي} \ar{سيارة} \ar{أبي}.
\end{quote}

\begin{quote}
\emph{aṣlaḥa al-mīkānīkī sayyārat abī.}
\end{quote}

\begin{quote}
The mechanic repaired my father's car.
\end{quote}

\begin{cousinweb}
\textbf{\ar{ميكانيكي}} (\emph{mīkānīkī}) is a modern loan, usually traced to French \emph{mécanique}, from Greek \emph{mēkhanē}, \textquotedblleft{}a device\textquotedblright{}, with the Arabic ending \emph{-ī}.
\end{cousinweb}

\begin{grammarlens}[title={What you've built}]
One more word for describing what you lived and felt.
\end{grammarlens}

\subsection*{Your turn}
\begin{itemize}
\raggedright
  \item \emph{Say it:} \emph{mīkānīkī}
  \item \emph{Say it:} \emph{mīkānīkī}, once more
  \item \emph{Say it:} \emph{aṣlaḥa al-mīkānīkī sayyārat abī.} — The mechanic repaired my father's car.
  \item \emph{Recall:} say the Arabic for to program, then the Arabic for a programmer, then say \emph{mīkānīkī} again
\end{itemize}

\begin{tcolorbox}[breakable,colback=teal!4,colframe=teal!35!black,title={Before you move on}]
What does \ar{ميكانيكي} mean? (\textquotedblleft{}A mechanic\textquotedblright{}.) Say it once more.
\end{tcolorbox}
\section[\texorpdfstring{\ar{مقاول} (muqāwil)}{مقاول (muqāwil)}]{\ar{مقاول} (muqāwil) — a building contractor}

\label{lesson:AR-C526-muqawil}



\begin{quote}
Before the new one: say the Arabic for a programmer, then the Arabic for a mechanic.
\end{quote}

\subsection*{You'll want to know: \ar{مقاول}}
\textbf{\ar{مقاول}} — \emph{muqāwil} — \textquotedblleft{}a building contractor\textquotedblright{}.

\begin{quote}
\ar{بنى} \ar{المقاول} \ar{بيتنا} \ar{الجديد}.
\end{quote}

\begin{quote}
\emph{banā al-muqāwil baytanā al-jadīd.}
\end{quote}

\begin{quote}
The contractor built our new house.
\end{quote}

\begin{cousinweb}
\textbf{\ar{مقاول}} (\emph{muqāwil}) is usually traced to the native Arabic root \emph{q-w-l}, \textquotedblleft{}to say\textquotedblright{}, the root of \emph{qāla}: Form III \emph{qāwala} is to talk terms with someone, and the \emph{muqāwil} is the one who takes on the work by agreement.
\end{cousinweb}

\begin{grammarlens}[title={What you've built}]
One more word for describing what you lived and felt.
\end{grammarlens}

\subsection*{Your turn}
\begin{itemize}
\raggedright
  \item \emph{Say it:} \emph{muqāwil}
  \item \emph{Say it:} \emph{muqāwil}, once more
  \item \emph{Say it:} \emph{banā al-muqāwil baytanā al-jadīd.} — The contractor built our new house.
  \item \emph{Recall:} say the Arabic for a programmer, then the Arabic for a mechanic, then say \emph{muqāwil} again
\end{itemize}

\begin{tcolorbox}[breakable,colback=teal!4,colframe=teal!35!black,title={Before you move on}]
What does \ar{مقاول} mean? (\textquotedblleft{}A building contractor\textquotedblright{}.) Say it once more.
\end{tcolorbox}
\section[\texorpdfstring{\ar{دهان} (dahhān)}{دهان (dahhān)}]{\ar{دهان} (dahhān) — a house painter}

\label{lesson:AR-C526-dahhan}



\begin{quote}
Before the new one: say the Arabic for a mechanic, then the Arabic for a building contractor.
\end{quote}

\subsection*{You'll want to know: \ar{دهان}}
\textbf{\ar{دهان}} — \emph{dahhān} — \textquotedblleft{}a house painter\textquotedblright{}.

\begin{quote}
\ar{دهن} \ar{الدهان} \ar{جدران} \ar{الغرفة}.
\end{quote}

\begin{quote}
\emph{dahana ad-dahhān judrān al-ghurfa.}
\end{quote}

\begin{quote}
The painter painted the walls of the room.
\end{quote}

The pattern faʿʿāl names a trade.

\begin{cousinweb}
\textbf{\ar{دهان}} (\emph{dahhān}) is built on the native Arabic root \emph{d-h-n}, \textquotedblleft{}oil, grease; to anoint\textquotedblright{} (Lane). Dahana, to paint a wall, which you know, is from the same root.
\end{cousinweb}

\begin{grammarlens}[title={What you've built}]
One more word for describing what you lived and felt.
\end{grammarlens}

\subsection*{Your turn}
\begin{itemize}
\raggedright
  \item \emph{Say it:} \emph{dahhān}
  \item \emph{Say it:} \emph{dahhān}, once more
  \item \emph{Say it:} \emph{dahana ad-dahhān judrān al-ghurfa.} — The painter painted the walls of the room.
  \item \emph{Recall:} say the Arabic for a mechanic, then the Arabic for a building contractor, then say \emph{dahhān} again
\end{itemize}

\begin{tcolorbox}[breakable,colback=teal!4,colframe=teal!35!black,title={Before you move on}]
What does \ar{دهان} mean? (\textquotedblleft{}A house painter\textquotedblright{}.) Say it once more.
\end{tcolorbox}
\section[\texorpdfstring{\ar{بناء} (bannāʾ)}{بناء (bannāʾ)}]{\ar{بناء} (bannāʾ) — a builder}

\label{lesson:AR-C526-banna}



\begin{quote}
Before the new one: say the Arabic for a building contractor, then the Arabic for a house painter.
\end{quote}

\subsection*{You'll want to know: \ar{بناء}}
\textbf{\ar{بناء}} — \emph{bannāʾ} — \textquotedblleft{}a builder\textquotedblright{}.

\begin{quote}
\ar{بنى} \ar{البناء} \ar{جدارا} \ar{حول} \ar{الحديقة}.
\end{quote}

\begin{quote}
\emph{banā al-bannāʾ jidāran ḥawla al-ḥadīqa.}
\end{quote}

\begin{quote}
The builder built a wall around the garden.
\end{quote}

Without the doubled n the same letters read bināʾ, a building.

\begin{cousinweb}
\textbf{\ar{بناء}} (\emph{bannāʾ}) is built on the native Arabic root \emph{b-n-y}, \textquotedblleft{}to build\textquotedblright{} (Lane). Banā, to build, which you know, is from the same root.
\end{cousinweb}

\begin{grammarlens}[title={What you've built}]
One more word for describing what you lived and felt.
\end{grammarlens}

\subsection*{Your turn}
\begin{itemize}
\raggedright
  \item \emph{Say it:} \emph{bannāʾ}
  \item \emph{Say it:} \emph{bannāʾ}, once more
  \item \emph{Say it:} \emph{banā al-bannāʾ jidāran ḥawla al-ḥadīqa.} — The builder built a wall around the garden.
  \item \emph{Recall:} say the Arabic for a building contractor, then the Arabic for a house painter, then say \emph{bannāʾ} again
\end{itemize}

\begin{tcolorbox}[breakable,colback=teal!4,colframe=teal!35!black,title={Before you move on}]
What does \ar{بناء} mean? (\textquotedblleft{}A builder\textquotedblright{}.) Say it once more.
\end{tcolorbox}
\section[\texorpdfstring{\ar{إسكافي} (iskāfī)}{إسكافي (iskāfī)}]{\ar{إسكافي} (iskāfī) — a cobbler}

\label{lesson:AR-C526-iskafi}



\begin{quote}
Before the new one: say the Arabic for a house painter, then the Arabic for a builder.
\end{quote}

\subsection*{You'll want to know: \ar{إسكافي}}
\textbf{\ar{إسكافي}} — \emph{iskāfī} — \textquotedblleft{}a cobbler\textquotedblright{}.

\begin{quote}
\ar{أصلح} \ar{الإسكافي} \ar{حذاء} \ar{أبي}.
\end{quote}

\begin{quote}
\emph{aṣlaḥa al-iskāfī ḥidhāʾ abī.}
\end{quote}

\begin{quote}
The cobbler repaired my father's shoe.
\end{quote}

\begin{cousinweb}
\textbf{\ar{إسكافي}} (\emph{iskāfī}) is an old Arabic word for a shoemaker or cobbler. Lane gives \emph{iskāf}, a maker of boots or shoes, under \emph{s-k-f}; the origin is not settled.
\end{cousinweb}

\begin{grammarlens}[title={What you've built}]
One more word for describing what you lived and felt.
\end{grammarlens}

\subsection*{Your turn}
\begin{itemize}
\raggedright
  \item \emph{Say it:} \emph{iskāfī}
  \item \emph{Say it:} \emph{iskāfī}, once more
  \item \emph{Say it:} \emph{aṣlaḥa al-iskāfī ḥidhāʾ abī.} — The cobbler repaired my father's shoe.
  \item \emph{Recall:} say the Arabic for a house painter, then the Arabic for a builder, then say \emph{iskāfī} again
\end{itemize}

\begin{tcolorbox}[breakable,colback=teal!4,colframe=teal!35!black,title={Before you move on}]
What does \ar{إسكافي} mean? (\textquotedblleft{}A cobbler\textquotedblright{}.) Say it once more.
\end{tcolorbox}
